\subsubsection{Reward B: Numeric Proximity}

\paragraph{Definition.}
Numeric proximity provides continuous partial credit based on the distance between the model's predicted answer and the gold answer.
If both are extractable as numbers:
\[
R_B(\text{pred}, \text{gold}) = \max\!\left(0,\; 1 - \frac{|\text{pred} - \text{gold}|}{|\text{gold}| + 1}\right)
\]
If the response does not contain a parseable number, $R_B = 0.0$.
The denominator $|\text{gold}| + 1$ normalizes the distance relative to the answer magnitude, and the $+1$ prevents division by zero when the gold answer is 0.

\paragraph{Rationale.}
Approximately 74\% of SFT baseline responses produce a correctly formatted answer that is numerically wrong.
Under binary correctness, all of these receive zero reward---a response off by 1 is treated identically to a response off by 10{,}000.
Numeric proximity provides graded signal proportional to how close the model's answer is, creating advantage variance among the 16 sampled responses even when none are exactly correct.
This is especially important for problems where the model consistently approaches the correct answer but makes small arithmetic errors.

\paragraph{Literature.}
The use of distance-based partial credit in reward shaping is grounded in the potential-based reward shaping framework of Ng et al.\ (1999), which shows that shaped rewards preserve optimal policies under certain conditions.
Unsloth's RL documentation recommends proximity-based rewards for math tasks to improve training stability when exact-match rates are low.

\paragraph{Complementarity with Reward A.}
Rewards A and B target non-overlapping failure populations.
A addresses the \textasciitilde11\% of responses lacking any parseable format---these receive $R_B = 0.0$ by definition since no number can be extracted.
B addresses the \textasciitilde74\% with correct format but wrong answer---these already satisfy $R_A = 1.0$.
There is no conflict: a response without format gets $(R_A, R_B) = (0, 0)$; a response with format but wrong answer gets $(1, \text{partial})$; a correct response gets $(1, 1)$.
Together, A and B provide non-zero signal for approximately 85\% of baseline failures.
